% ============================================================
%  NOMENCLATURA
% ============================================================
\section*{Nomenclatura}

% ── A. Símbolos latinos ──────────────────────────────────────
\subsection*{A. Símbolos latinos}

\begin{longtable}{@{} l p{0.65\textwidth} r @{}}
\toprule
\textbf{Símbolo} & \textbf{Descripción} & \textbf{Unidad} \\
\midrule

$A$
  & Matriz de estado del modelo en espacio de estados (Ecs.~\ref{eq:statespace_state}--\ref{eq:statespace_output})
  & -- \\

$A(q)$
  & Polinomio autorregresivo del modelo ARX/ARMAX (Ecs.~\ref{eq:arx_general}--\ref{eq:arx_A})
  & -- \\

$a_i$
  & Coeficientes del polinomio $A(q)$, $i=1,\ldots,n_a$
  & -- \\

$B$
  & Matriz de entrada del modelo en espacio de estados (Ec.~\ref{eq:statespace_state})
  & -- \\

$B(q)$
  & Polinomio de entrada del modelo ARX/ARMAX (Ecs.~\ref{eq:arx_general},~\ref{eq:arx_B})
  & -- \\

$b_i$
  & Coeficientes del polinomio $B(q)$, $i=1,\ldots,n_b$
  & -- \\

$C$
  & Matriz de salida del modelo en espacio de estados (Ec.~\ref{eq:statespace_output})
  & -- \\

$C$
  & Número de componentes químicos en el modelo MESH de destilación (Ec.~\ref{eq:num_ec})
  & -- \\

$C(q)$
  & Polinomio de ruido del modelo ARMAX (Ec.~\ref{eq:armax_C})
  & -- \\

$c_i$
  & Coeficientes del polinomio $C(q)$, $i=1,\ldots,n_c$
  & -- \\

$e(k)$
  & Término de ruido blanco / perturbación no modelada en ARX (tiempo discreto)
  & -- \\

$\bar{e}(k)$
  & Señal de ruido en tiempo discreto en la formulación ARX/ARMAX (Ec.~\ref{eq:arx_general})
  & -- \\

$e(t)$
  & Error de seguimiento de control (tiempo continuo), utilizado en IAE, ISE, ITAE, ITSE
  & $^{\circ}$C \\

$e_{rel}(t)$
  & Error relativo de seguimiento entre la variable controlada y el setpoint (Ec.~\ref{eq:ErrorRelativo})
  & -- \\

$er_1,\,er_2$
  & Errores relativos en estado estacionario tras la primera y segunda perturbación (Tabla~\ref{tab:ControlIntegralesP1}--\ref{tab:ControlIntegralesP2})
  & -- \\

$i$
  & Índice de paso de predicción dentro del horizonte MPC ($i=1,\ldots,p$)
  & -- \\

$j$
  & Índice de variable controlada o variable manipulada
  & -- \\

$J(z_k)$
  & Función de costo (objetivo) total del MPC (Ec.~\ref{eq:OptimizationProblem})
  & -- \\

$J_y(z_k)$
  & Término de seguimiento de salida de la función de costo MPC (Ec.~\ref{eq:Jy})
  & -- \\

$J_u(z_k)$
  & Término de penalización de variables manipuladas de la función de costo MPC (Ec.~\ref{eq:Ju})
  & -- \\

$J_{\Delta u}(z_k)$
  & Término de penalización de tasa de cambio de las acciones de control (Ec.~\ref{eq:Jdeltau})
  & -- \\

$J_{\varepsilon}(z_k)$
  & Término de penalización de la variable de holgura (Ec.~\ref{eq:Jepsilon})
  & -- \\

$k$
  & Índice de paso de tiempo discreto actual
  & -- \\

$n$
  & Índice de muestra utilizado en las sumatorias de RECM y $R^2$ (Ecs.~\ref{eq:RECM}--\ref{eq:R2})
  & -- \\

$n_a$
  & Orden del polinomio autorregresivo $A(q)$ (número de polos)
  & -- \\

$n_b$
  & Orden del polinomio de entrada $B(q)$ (número de ceros)
  & -- \\

$n_c$
  & Orden del polinomio de ruido $C(q)$ en el modelo ARMAX
  & -- \\

$n_k$
  & Retardo de entrada del modelo ARX/ARMAX (número de muestras) (Ec.~\ref{eq:arx_general})
  & -- \\

$n_u$
  & Número de variables manipuladas en la formulación MPC
  & -- \\

$n_y$
  & Número de variables controladas en la formulación MPC
  & -- \\

$N$
  & Número de muestras de datos utilizadas en los cálculos de RECM y $R^2$ (Ecs.~\ref{eq:RECM}--\ref{eq:R2})
  & -- \\

$N$
  & Número de etapas teóricas de la columna de destilación (modelo MESH, Ec.~\ref{eq:num_ec})
  & -- \\

$N_{ec}$
  & Número de ecuaciones algebraicas que describen la columna en un instante dado (Ec.~\ref{eq:num_ec})
  & -- \\

$N_{et}$
  & Número total de ecuaciones resueltas durante la simulación dinámica (Ec.~\ref{eq:num_et})
  & -- \\

$N_i$
  & Número de intervalos de tiempo en la simulación dinámica (Ec.~\ref{eq:num_et})
  & -- \\

$p$
  & Horizonte de predicción del MPC (Ecs.~\ref{eq:Vectorzk}--\ref{eq:Jdeltau})
  & pasos \\

$P_1$
  & Temperatura de la variable controlada tras la primera perturbación (Tabla~\ref{tab:ControlIntegralesP1}--\ref{tab:ControlIntegralesP2})
  & $^{\circ}$C \\

$P_2$
  & Temperatura de la variable controlada tras la segunda perturbación (Tabla~\ref{tab:ControlIntegralesP1}--\ref{tab:ControlIntegralesP2})
  & $^{\circ}$C \\

$q$
  & Operador de avance (adelanto)
  & -- \\

$q^{-1}$
  & Operador de retardo unitario
  & -- \\

$Q_{R}$
  & Calor del rehervidor (genérico)
  & kW \\

$Q_{R,T\text{-}101}$
  & Calor del rehervidor de la columna extractiva T-101 (variable manipulada)
  & kW \\

$Q_{R,T\text{-}102}$
  & Calor del rehervidor de la columna de recuperación de solvente T-102 (variable manipulada)
  & kW \\

$R^2$
  & Coeficiente de determinación (métrica de ajuste, Ec.~\ref{eq:R2})
  & -- \\

$r_j(k{+}i|k)$
  & Referencia predicha (setpoint) para la salida $j$ en el paso de predicción $i$ (Ec.~\ref{eq:Jy})
  & -- \\

$RECM$
  & Raíz del Error Cuadrático Medio (Ec.~\ref{eq:RECM})
  & $^{\circ}$C \\

$RR_{T\text{-}101}$
  & Razón de reflujo de la columna T-101 (variable manipulada)
  & -- \\

$RR_{T\text{-}102}$
  & Razón de reflujo de la columna T-102 (variable manipulada)
  & -- \\

$s_j^u$
  & Factor de escala para la variable manipulada $j$ en el costo MPC (Ecs.~\ref{eq:Ju}--\ref{eq:Jdeltau})
  & -- \\

$s_j^y$
  & Factor de escala para la variable de salida $j$ en el costo MPC (Ec.~\ref{eq:Jy})
  & -- \\

$SP$
  & Temperatura de referencia (setpoint) (Tabla~\ref{tab:ControlIntegralesP1}--\ref{tab:ControlIntegralesP2})
  & $^{\circ}$C \\

$t$
  & Tiempo continuo; también variable de integración en IAE, ISE, ITAE, ITSE
  & h \\

$t_1$
  & Tiempo de estabilización tras la primera perturbación (Tabla~\ref{tab:ControlIntegralesP1}--\ref{tab:ControlIntegralesP2})
  & h \\

$t_2$
  & Tiempo de estabilización tras la segunda perturbación (Tabla~\ref{tab:ControlIntegralesP1}--\ref{tab:ControlIntegralesP2})
  & h \\

$T$
  & Temperatura (genérica)
  & $^{\circ}$C \\

$T_{36,T\text{-}101}$
  & Temperatura en la etapa 36 de la columna T-101 — variable controlada de T-101
  & $^{\circ}$C \\

$T_{8,T\text{-}102}$
  & Temperatura en la etapa 8 de la columna T-102 — variable controlada de T-102
  & $^{\circ}$C \\

$u$
  & Variable manipulada (genérica)
  & -- \\

$u(k)$
  & Señal de variable manipulada en tiempo discreto
  & -- \\

$\bar{u}(k)$
  & Señal de entrada en tiempo discreto en la formulación ARX/ARMAX (Ec.~\ref{eq:arx_general})
  & -- \\

$u(t)$
  & Vector de entradas en tiempo continuo (espacio de estados, Ec.~\ref{eq:statespace_state})
  & -- \\

$u_j(k{+}i|k)$
  & Variable manipulada $j$ predicha en el paso $i$ (Ecs.~\ref{eq:Ju}--\ref{eq:Jdeltau})
  & -- \\

$u_{j,target}(k{+}i|k)$
  & Valor objetivo (estado estable deseado) para la variable manipulada $j$ (Ec.~\ref{eq:Ju})
  & -- \\

$w_{i,j}^{y}$
  & Peso de sintonización MPC para la salida $j$ en el paso $i$ (Ec.~\ref{eq:Jy})
  & -- \\

$w_{i,j}^{u}$
  & Peso de sintonización MPC para la variable manipulada $j$ en el paso $i$ (Ec.~\ref{eq:Ju})
  & -- \\

$w_{i,j}^{\Delta u}$
  & Peso de sintonización MPC para la tasa de cambio de la variable manipulada $j$ (Ec.~\ref{eq:Jdeltau})
  & -- \\

$x(t)$
  & Vector de estados en tiempo continuo (Ec.~\ref{eq:statespace_state})
  & -- \\

$\dot{x}(t)$
  & Derivada temporal del vector de estados (Ec.~\ref{eq:statespace_state})
  & -- \\

$X$
  & Variable de proceso genérica utilizada en la normalización por desviación (Ec.~\ref{eq:Desviacion})
  & -- \\

$X_{desv}$
  & Desviación relativa de la variable $X$ respecto a su valor nominal de operación (Ec.~\ref{eq:Desviacion})
  & -- \\

$X_{nominal}$
  & Valor nominal (estado estacionario) de la variable $X$ (Ec.~\ref{eq:Desviacion})
  & -- \\

$y$
  & Variable de salida controlada (genérica)
  & -- \\

$y(k)$
  & Señal de salida del sistema en tiempo discreto (formulación ARX/ARMAX)
  & -- \\

$\bar{y}(k)$
  & Señal de salida en tiempo discreto en la formulación ARX/ARMAX (Ec.~\ref{eq:arx_general})
  & -- \\

$y(t)$
  & Vector de salidas en tiempo continuo (espacio de estados, Ec.~\ref{eq:statespace_output})
  & -- \\

$y_j(k{+}i|k)$
  & Salida controlada $j$ predicha en el paso $i$ (Ec.~\ref{eq:Jy})
  & -- \\

$y_n$
  & Salida observada (modelo riguroso) en la muestra $n$ (Ecs.~\ref{eq:RECM}--\ref{eq:R2})
  & -- \\

$\hat{y}_n$
  & Salida predicha por el modelo subrogado en la muestra $n$ (Ecs.~\ref{eq:RECM}--\ref{eq:R2})
  & -- \\

$\bar{y}$
  & Media de las muestras de salida observadas; utilizada en el denominador de $R^2$ (Ec.~\ref{eq:R2})
  & -- \\

$Y(t)$
  & Salida del proceso controlado en tiempo continuo (Ec.~\ref{eq:ErrorRelativo})
  & $^{\circ}$C \\

$Y_{ref}(t)$
  & Valor de referencia (setpoint) de la salida del proceso (Ec.~\ref{eq:ErrorRelativo})
  & $^{\circ}$C \\

$z_k$
  & Vector de decisión MPC — secuencia de acciones de control futuras y variable de holgura (Ec.~\ref{eq:Vectorzk})
  & -- \\

$z_k^T$
  & Transpuesta del vector de decisión MPC (Ec.~\ref{eq:Vectorzk})
  & -- \\

\bottomrule
\end{longtable}

\bigskip

% ── B. Símbolos griegos ──────────────────────────────────────
\subsection*{B. Símbolos griegos}

\begin{tabularx}{\textwidth}{@{} l X r @{}}
\toprule
\textbf{Símbolo} & \textbf{Descripción} & \textbf{Unidad} \\
\midrule

$\Delta$
  & Operador de diferencia finita / incremento (p.ej.\ $\Delta u$: cambio en la variable manipulada)
  & -- \\

$\varepsilon_k$
  & Variable de holgura para la relajación de restricciones blandas en el paso $k$
    (escrita como $\epsilon_k$ en la Ec.~\ref{eq:Vectorzk} y como $\varepsilon_k$ en la Ec.~\ref{eq:Jepsilon};
    ambas se refieren a la misma variable)
  & -- \\

$\rho_{\varepsilon}$
  & Coeficiente de penalización sobre la variable de holgura en el costo MPC (Ec.~\ref{eq:Jepsilon})
  & -- \\

\bottomrule
\end{tabularx}

\bigskip

% ── C. Subíndices ────────────────────────────────────────────
\subsection*{C. Subíndices}

\begin{tabularx}{\textwidth}{@{} l X @{}}
\toprule
\textbf{Subíndice} & \textbf{Descripción} \\
\midrule
$36$                & Etapa 36 de la columna T-101 (ubicación de medición de temperatura) \\
$8$                 & Etapa 8 de la columna T-102 (ubicación de medición de temperatura) \\
$a$                 & Orden del polinomio autorregresivo \\
$b$                 & Orden del polinomio de entrada \\
$c$                 & Orden del polinomio de ruido \\
$ec$                & Conteo de ecuaciones algebraicas (modelo MESH) \\
$et$                & Conteo total de ecuaciones en la simulación dinámica \\
$i$                 & Índice de paso de predicción dentro del horizonte MPC \\
$j$                 & Índice de variable controlada o variable manipulada \\
$k$                 & Paso de tiempo discreto actual \\
$n$                 & Índice de muestra en las métricas de ajuste estadístico (RECM, $R^2$) \\
$R$                 & Rehervidor (utilizado en $Q_{R}$, $Q_{R,T\text{-}101}$, etc.) \\
$ref$               & Valor de referencia (setpoint) \\
$T\text{-}101$      & Identificador de la columna de destilación extractiva \\
$T\text{-}102$      & Identificador de la columna de recuperación de solvente \\
\bottomrule
\end{tabularx}

\bigskip

% ── E. Abreviaturas ──────────────────────────────────────────
\subsection*{E. Abreviaturas}

\begin{tabularx}{\textwidth}{@{} l X @{}}
\toprule
\textbf{Abreviatura} & \textbf{Descripción} \\
\midrule
ARMAX       & AutoRegressive Moving Average with eXogenous inputs (Autorregresivo de media móvil con entradas exógenas) \\
ARX         & AutoRegressive with eXogenous inputs (Autorregresivo con entradas exógenas) \\
NLARX       & NonLinear AutoRegressive with eXogenous inputs (Autorregresivo no lineal con entradas exógenas) \\
CPU         & Central Processing Unit (Unidad Central de Procesamiento) \\
DMSO        & Dimetilsulfóxido (solvente de destilación extractiva) \\
FOPDT       & First-Order Plus Dead Time (Primer orden con tiempo muerto) \\
IAE         & Integral of Absolute Error — Integral del valor absoluto del error (Ec.~\ref{eq:IAE}) \\
ISE         & Integral of Squared Error — Integral del error cuadrático (Ec.~\ref{eq:ISE}) \\
ITAE        & Integral of Time-weighted Absolute Error — Integral del valor absoluto del error ponderado por el tiempo (Ec.~\ref{eq:ITAE}) \\
ITSE        & Integral of Time-weighted Squared Error — Integral del error cuadrático ponderado por el tiempo (Ec.~\ref{eq:ITSE}) \\
MESH        & Mass, Equilibrium, Summation and Heat (marco de modelado de columnas de destilación) \\
MIMO        & Multiple Inputs, Multiple Outputs (múltiples entradas, múltiples salidas) \\
MISO        & Multiple Inputs, Single Output (múltiples entradas, una salida) \\
MPC         & Model Predictive Control — Control Predictivo Basado en Modelo \\
PFD         & Process Flow Diagram — Diagrama de Flujo del Proceso \\
PI          & Controlador Proporcional-Integral \\
PID         & Controlador Proporcional-Integral-Derivativo \\
QP          & Quadratic Programming — Programación Cuadrática \\
RAM         & Random Access Memory — Memoria de Acceso Aleatorio \\
RECM        & Raíz del Error Cuadrático Medio (Ec.~\ref{eq:RECM}) \\
SISO        & Single Input, Single Output (una entrada, una salida) \\
SS          & State Space — Espacio de Estados \\
\bottomrule
\end{tabularx}